\section{標準補完接線と対称性障害の比較資料}
\label{sec:completion-notes-v38}

\subsection{継続に必要な確定事項}

$m=3$, $j=r=1$ の量子VK corner $B$ と標準六項代数 $\Lambda$ について、
補完接線 $\ell=\operatorname{im}(HH^2(\Lambda)\to HH^2(B))$ は
\[
 \ell\subset K=\ker[\delta,-],\qquad
 \ell\cap\operatorname{im}[\delta,-]=0,
 \qquad \delta=\log(1+N)
\]
を満たす。従って、v37で残っていた「唯一の補完接線は、五次元の
$K\cap\operatorname{im}[\delta,-]$ に入るか」という問題の答えは否定である。
一次方向を失わずに、二次で外的導分を失う。

非零の補完初速度を持つ任意の形式弧 $\mathcal O=\mathbf C[[t]]$ に沿って、
対称性保持のFittingイデアルは $(t^2)$ へ引き戻される。
従って最小相対複体は
\[
 \mathcal O[0]\oplus
 (\mathcal O\xrightarrow{t^2}\mathcal O)_{1,2}
 \oplus\mathcal O^{34}[-2]\oplus\mathcal O^{39}[-3]
\]
となる。相対 $HH^2$ の捩れは正確に $\mathcal O/(t^2)$ であり、
一般ファイバーの $HH^1$ は零となる。
この結論は、二次以上の積補正や、滑らかな三十五次元hull内の弧の選択によらない。

初速度だけで弧を分類すると、$B_{1,1}$ では対称性消失次数は
$u\notin K$ で一、$u\in K\setminus\operatorname{im}L$ で二に固定される。
$u\in K\cap\operatorname{im}L$ では、任意の整数 $n\ge2$ と無限大が実現する。
後者は、滑らかな五次元部分 $Y$ 上の弧 $c(t)$ に
$t^n v$（$v\notin K$）を加えることで有限次数 $n$、$c(t)$ 自身で無限大を得る。
従って「二次対称性障害が零」という条件は、二次補正を適切に選べることを意味し、
任意の二次補正が対称性を保つことを意味しない。
$B_{0,3}$ の四次元の不変接空間についても、同じ次数の全範囲が実現する。

\subsection{再使用できるOre構成}

本計算では $\alpha=-1$, $s=1$ の表示で一次の外側変形を作り、
有限Zhang同型で v37 の $\alpha=0$ 表示へ戻した。
三変数代数は、既存の $C_{\gamma,s}$ の $s=1$, $\gamma^2=0$ への特殊化である。
$ux=x\sigma(u)$ を満たす正規元 $x$ を用いて
\[
 \mathcal S=C[w;\sigma^2,D_H],\qquad
 D_H(u)=x^{-2}(uH-Hu)
\]
とする。必要なのは $[C,H]\subset x^2C$ だけであり、
このとき $x^2w+H$ が中心元となる。
この補題は任意の正整数 $k$ に対し、
\[
 [C,H]\subset x^kC,\qquad
 D_H=x^{-k}[-,H],\qquad
 C[w;\sigma^k,D_H],\qquad J=x^kw+H
\]
へそのまま一般化できる。
$r>0$ の他の共鳴で $k=r+1$ を使う場合に再利用できる。
端点 $r=0$ の構成は $k=1$ に当たり、商 $C/(x)$ の可換性だけで
任意の $H$ を許せた理由もこの形で統一される。

今回の $H_0$ はCGG中心元
\[
 J=x^2w+q\left(xyz-\frac23y^3\right)
\]
のPBW表示から得た。一次補正 $H_1$ は、次数五のPBW元を未知数として
$[x,H],[y,H],[Z,H]$ の $x$ 次数零・一の係数を消す線形方程式を解いて得た。
未知数二十一、制約行三十九、係数行列と拡大行列の階数は共に十である。
自由変数を零にした解は公表稿に明示した五項の $H_1$ である。
次数三の $G$ に $H=QG$ という因子分解を課すと、
対応する行列の階数は六、拡大行列は七となり一次ですでに解がなくなる。
従って、$H_\varepsilon$ が二次の中心元 $Q_\varepsilon$ の倍数であるという
余分な制約を課すと、存在する補完接線を誤って排除する。
この計算は構成法の選択に関わるため保持する。

\subsection{計算規約と保存データ}

経路順の積を用い、cornerの対象は $(0,1,3,4)$ とする。
$\alpha=0$ から $\alpha=-1$ への有限同型は、$i\to k$ の射上で
$F_{ik}=T^{-k}$ である。
さらに、PBW元 $Z$ を $z+y$ へ送る固定した基底変換を用いる。
生成した一次cochain $U$ は $1584$ 成分を持つ。
v37の核基底における第六座標は $-1571/60$ であり、
余核の二次障害の第二座標は
\[
 \frac{432000}{2468041}\left(-\frac{1571}{60}\right)^2=120.
\]
この一つの非零係数だけでも、補完接線が二次対称性障害を持つことが分かる。

\texttt{verify\_vk\_completion\_symmetry\_v38.py} は通常実行時に書込みを行わず、
PBWの重なり、正規性、自己同型、全六対象における特殊ファイバーの積と一次結合律、
cornerのcocycle条件、核・余核座標、二種類の二次障害計算を行う。
証明書 \texttt{vk\_completion\_symmetry\_v38\_certificate.json} は
$U$、$d\eta=L_\delta U$ の解、$dV=U\circ U$ の解、
$L_\delta V+L_\eta U$ のコホモロジー座標を保持する。
比較先のv37証明書のSHA--256も検査する。

\subsection{全局所環の判定と残る障害の比較}

補完接線上の二次障害が非零であることだけから、全対称性保持部分の
完備局所環や非被約性を推論してはならない。
今回の全環の決定には、五変数の厳密な $\mathbf G_a$ 族に
平方零横断方向を加え、全ての結合則と導分条件を多項式恒等式として確認する議論が必要であった。
その結果、$A=\mathbf C[[u_1,\ldots,u_5]]$ に対し全環は
$A[[\gamma]]/(\gamma^2)$ と決まり、以前の有限表示における $b,J$ は共に零となる。
純粋部分への制限と適合する座標では、標準六項補完像は $u_1=\cdots=u_5=0$、
被約部分 $Y$ は $\gamma=0$ である。
この埋込みを含む同時正規形までが確定した範囲である。

一方、外側の結合性障害は $HH^3(\Lambda)$、内側の対称性障害は
$\operatorname{coker}(L|HH^2(B))$ に値を持つ。
両者の自然な比較は、後続の導分付き障害複体によって解決した。
外側の結合性障害に導分の一次補正を加えた類は、全充満部分圏への制限で自然に写る。
内側の結合性誤差を消すと、対称性作用の余核に一意の類が残る。
標準六項では一次元の一次障害線上に自然な非零線形写像を得る。
七次元の $HH^3$ 全体への比較写像や、一般重みの高次比較はこの主張に含めない。
